\section{Introduction}
For many real-world application domains such as robotics or virtual assistants, collecting large amounts of data for training an agent can be time consuming, expensive, or even dangerous.
Hence, the ability to learn new tasks from a handful of demonstrations is crucial for enabling a wider range of use cases. 
However, current deep learning agents often struggle to learn new tasks from a limited number of demonstrations.

Prior work attempts to address this problem using meta-learning~\citep{schmidhuber1987evolutionary,duan2016rl,finn2017model,pong2022offline, mitchell2021offline,beck2023survey} but these methods tend to be difficult to use in practice and require more than a handful of demonstrations or extensive fine-tuning. 
In contrast, large transformers trained on vast amounts of data can learn new tasks from only a few examples without any parameter updates~\citep{gpt3, scaling-laws, olsson2022incontext, zipfian_icl}. This emergent phenomenon is called \textit{few-shot or in-context learning (ICL)}, and is achieved by simply conditioning the model's outputs on a context containing a few examples for solving the task~\citep{gpt3, Ganguli_2022, wei2022emergent}. 

While in-context learning has been observed in multiple domains from language~\citep{gpt3} to vision~\citep{zipfian_icl}, it has not yet been extensively studied in sequential decision making settings. The sequential decision making problem poses additional challenges that do not appear in the supervised or self-supervised settings. For example, this setting tends to have a lower tolerance for errors since the environment's stochasticity or the agent's actions can lead to unseen, and sometimes unrecoverable, states. 

In this paper, we study in-context learning of sequential decision making tasks. To test our approach, we consider a challenging setting where the train and the test task distributions are disjoint with different states, actions, dynamics, and reward functions. For example, on Procgen we train on 12 of the games and test on the remaining 4. To emphasize how different the train and test tasks are, one of the training tasks is Bigfish where the must eat as many smaller fish as possible while avoiding larger fish than itself and it grows in size as it eats more. In contrast, one of our test tasks is Ninja, where the agent must jump across ledges while avoiding bombs or destroying them by tossing starsin order to collect the mushroom at the end of the platform. \Cref{fig:train_test} illustrates some of our train and test task, emphasizing how different they are. To our knowledge, no prior work has demonstrated the ability to generalize to new MiniHack or Procgen tasks using only a handful of demonstrations and no weight updates.


We first focus on the data distributional properties required to enable in-context learning of sequential decision making tasks. 
Our key finding is that in contrast to (self-)supervised learning where the context can simply contain a few different examples (or predictions), in sequential decision making it is crucial for the context to contain full/partial trajectories (or sequences of predictions) to cover the potentially wide range of states the agent may find itself in at deployment. 
Translating this insight into a dataset construction pipeline, we demonstrate that we can enable in-context learning of \textit{unseen} tasks on both the MiniHack and Procgen benchmarks from only a handful of demonstrations.
We additionally perform an extensive study of other design choices such as the model size, dataset size, trajectory burstiness, environment stochasticity, and task diversity. Our experiments suggest that larger model and dataset sizes, as well as more trajectory burstiness, environment stochasticity, and task diversity, all lead to better in-context learning when using transformers for sequential decision making tasks. \textbf{Our work is first to show that generalization to new Procgen and MiniHack tasks is possible from just a handful of expert demonstrations and no weight updates using the transformer's in-context learning ability.} This goes beyond the well-studied generalization to new levels which are merely procedurally generated instances of the training task rather than fundamentally new tasks~\citep{cobbe2020leveraging, igl2019generalization, laskin2020reinforcement, Raileanu2020AutomaticDA, raileanu2021decoupling}. 

\begin{figure}[t]
    \centering
    \includegraphics[width=\columnwidth]{figures/Procgen_games.png}
    \caption{\textbf{Illustration of Train and Test Tasks.} (\textit{Left}) A collection of procedurally generated Procgen levels from the \texttt{Fruitbot} task, demonstrating the complexity and diversity inherent in the environment's design. (\textit{Middle}) Tasks used for training. (\textit{Right}) Tasks used for testing. Note that the test tasks are entirely distinct from the training tasks, and each of them is procedurally generated, consisting of multiple levels.}
    \label{fig:train_test}
\end{figure}
